% Part: normal-modal-logic
% Chapter: filtrations
% Section: euclidean-filtrations

\documentclass[../../../include/open-logic-section]{subfiles}

\begin{document}

\olfileid{nml}{fil}{euc}

\olsection{युक्लिडी प्रतिमानांची निस्यंदने}

\olref[acc]{sec} मधील पद्धत युक्लिडी, किंवा
सर्वत्र उत्तरवर्ती आणि युक्लिडी, प्रतिमानांसाठी
चालत नाही. \olref{fig:ser-eucl} च्या वरच्या भागातील
प्रतिमान पाहा; ते युक्लिडी आणि सर्वत्र उत्तरवर्ती आहे.
$\Gamma = \{p, \Box p \}$ असू द्या. $\Gamma$ मधून
निस्यंदन घेतल्यावर $[w_1] = [w_3]$, कारण $w_1$
आणि $w_3$ हीच दोन जगे~$\Gamma$ वरील सत्यतेबाबत
एकमत आहेत. कोणत्याही निस्यंदनात~$\mModel{M}$
मधून वारशाने आलेले बाणही असतील; ते
\olref{fig:ser-eucl2} मध्ये दाखवले आहेत. ते प्रतिमान
युक्लिडी नाही. शिवाय ते युक्लिडी करण्यासाठी
त्यात आणखी बाण घालता येत नाहीत. $[w_2]$
आणि $[w_4]$ यांच्यामध्ये दोन्ही दिशांनी बाण
घालावे लागतील; मग $[w_2]$ आणि~$[w_5]$
यांच्यामध्येही. पण~$w_2$ येथे $\Box p$ सत्य
असायला हवे, तर~$w_5$ येथे $p$ असत्य आहे.

\begin{figure}[htpb]
  \centering
  \begin{tikzpicture}[modal]
    \node[world] (w1)
          [label={left: \mFalse{p}},
            label={below: $\mSat{{}}{\Box p}$}] {$w_1$};
    \node[world] (w2)
          [label={right: \mTrue{p}},
            label = {below: $\mSat{{}}{\Box p}$},
            right=of w1] {$w_2$};
    \draw[reflexive above] (w2) to (w2);
    \draw[->] (w1) to (w2);
    \node[world] (w3)
          [label={left: \mFalse{p}},
            label={below: $\mSat{{}}{\Box p}$},
            below=of w1] {$w_3$};
    \node[world] (w4)
          [label={right:\mTrue{p}},
            label={below: $\mSat/{{}}{\Box p}$},
            right=of w3] {$w_4$};
    \draw[->] (w3) to (w4);
    \draw[reflexive above] (w4) to (w4);
    \node[world] (w5)
          [label={right: \mFalse{p}},
            label=below:$\mSat/{{}}{\Box p}$,
            right=of w4] {$w_5$};
    \draw[->, bend left] (w4) to (w5);
    \draw[->, bend left] (w5) to (w4);
    \draw[reflexive above] (w5) to (w5);
  \end{tikzpicture}
  \caption{सर्वत्र उत्तरवर्ती आणि युक्लिडी प्रतिमान.}
  \ollabel{fig:ser-eucl}
\end{figure}

\begin{figure}[ht]
  \centering
  \begin{tikzpicture}[modal,every text node part/.style={align=center}]
    \node[world] (w1)
          [label={left: \mFalse{p}},
            label={right:$[w_1]=[w_3]$},
            label={below: $\mSat{{}}{\Box p}$}] {$[w_1]$};
    \node[world] (w2)
         [label={right: \mTrue{p}},
           label = {below: $\mSat{{}}{\Box p}$},
           right=of w1, yshift=1.5cm] {$[w_2]$};
    \draw[reflexive above] (w2) to (w2);
    \draw[->] (w1) to (w2);
    \node[world] (w4)
         [label={right:\mTrue{p}},
           label={below: $\mSat/{{}}{\Box p}$},
           right=of w1,yshift=-1.5cm] {$[w_4]$};
    \draw[->] (w1) to (w4);
    \draw[reflexive above] (w4) to (w4);
    \node[world] (w5)
         [label={right: \mFalse{p}},
             label=below:$\mSat/{{}}{\Box p}$,
             right=of w4] {$[w_5]$};
    \draw[->, bend left] (w4) to (w5);
    \draw[->, bend left] (w5) to (w4);
    \draw[reflexive above] (w5) to (w5);
  \end{tikzpicture}
  \caption{\olref{fig:ser-eucl} मधील प्रतिमानाचे निस्यंदन.}
  \ollabel{fig:ser-eucl2}
\end{figure}

विशेषतः युक्लिडी निस्यंदन मिळवण्यासाठी
उप-!!{formula}s खाली बंद असलेल्या मनमानी $\Gamma$
मधून निस्यंदन घेणे पुरेसे नाही. त्याऐवजी
\emph{मोडलदृष्ट्या बंद} असलेले संच~$\Gamma$
घ्यावे लागतात (\olref[pre]{defn:modallyclosed} पाहा).
असे वाक्यांचे संच अनंत असतात; म्हणून त्यांच्यापासून
संबंधित प्रणालीचा सांत प्रतिमान गुणधर्म किंवा
निर्णेयता तात्काळ मिळत नाही.

\begin{thm}\ollabel{thm:modal-closed-filt}
  $\Gamma$ मोडलदृष्ट्या बंद असू द्या,
  $\mModel{M}=\tuple{W,R,V}$ हे प्रतिमान असू द्या,
  आणि $\mModel{M^*} = \tuple{W^*,R^*,V^*}$ हे
  $\mModel{M}$ चे सर्वांत स्थूल निस्यंदन असू द्या.
  \begin{enumerate}
  \item $\mModel{M}$ सममित असेल, तर $\mModel{M^*}$ सुद्धा सममित आहे.
  \item $\mModel{M}$ संक्रमक असेल, तर $\mModel{M^*}$ सुद्धा संक्रमक आहे.
  \item $\mModel{M}$ युक्लिडी असेल, तर $\mModel{M^*}$ सुद्धा युक्लिडी आहे.
  \end{enumerate}
\end{thm}

\begin{proof}
\begin{enumerate}
  \item सराव. \Ax{B} आणि $\Ax{B_\Diamond}$ ही दोन्ही
    सर्व सममित प्रतिमानांत वैध आहेत, हे वापरा.
  \item $\mModel{M^*}$ हे सर्वांत स्थूल निस्यंदन असेल,
    तर व्याख्येनुसार $R^*[u][v]$ तेव्हा आणि तेव्हाच
    $C_1(u,v)$. संक्रमकतेसाठी $C_1(u,v)$ आणि
    $C_1(v,w)$ असे समजा; $C_1(u,w)$ दाखवायचे आहे.
    \iftag{prvBox}{$\mSat{M}{\Box !A}[u]$ असे समजा.
      सर्व संक्रमक प्रतिमानांत \Ax{4} वैध असल्याने
      $\mSat{M}{\Box\Box!A}[u]$. मोडल बंदतेवरून
      $\Box\Box!A \in \Gamma$; म्हणून $C_1(u,v)$ वरून
      $\mSat{M}{\Box!A}[v]$, आणि $C_1(v,w)$ वरून
      $\mSat{M}{!A}[w]$. }{}\iftag{prvDiamond}{
      $\mSat{M}{!A}[w]$ असे समजा. मोडल बंदतेवरून
      $\Diamond !A \in \Gamma$; म्हणून $C_1(v,w)$ वरून
      $\mSat{M}{\Diamond !A}[v]$. तसेच मोडल बंदतेवरून
      $\Diamond\Diamond!A \in \Gamma$; म्हणून $C_1(u,v)$ वरून
      $\mSat{M}{\Diamond\Diamond !A}[u]$.
      सर्व संक्रमक प्रतिमानांत $\Ax{4}_\Diamond$
      वैध असल्याने $\mSat{M}{\Diamond!A}[u]$.}{}
  \item सराव. \Ax{5} आणि $\Ax{5_\Diamond}$ ही दोन्ही
    सर्व युक्लिडी प्रतिमानांत वैध आहेत, हे वापरा.
\end{enumerate}
\end{proof}

\begin{prob}
  \olref[nml][fil][euc]{thm:modal-closed-filt} ची सिद्धता पूर्ण करा.
\end{prob}

\end{document}
